\documentclass[11pt,a4paper]{article}

\usepackage[utf8]{inputenc}
\usepackage[T1]{fontenc}
\usepackage[margin=2.5cm]{geometry}
\usepackage{amsmath,amssymb,bm}
\usepackage{booktabs}
\usepackage{array}
\usepackage{enumerate}
\usepackage{xcolor}
\usepackage{hyperref}
\usepackage{tikz}

\hypersetup{colorlinks=true, linkcolor=blue, urlcolor=blue}

\newcommand{\points}[1]{\textbf{(#1\,\%)}}
\newcommand{\subpoints}[1]{\textit{(#1\,\%)}}
\newcommand{\sol}{\noindent\textbf{Solution}}

\title{TMA4268 V2026 Mock Exam 6 --- Solution Proposal\\[0.3em]
\large}
\author{Compiled for Anders Bekkevard\\[0.2em]
\small Companion to \texttt{mock-exam-6.tex} (same directory).}
\date{Mock for: May 18, 2026}

\begin{document}
\maketitle

\noindent\textit{This document is a worked-solution proposal in the style of the official
Stefanie/Sara solutions to the 2024 and 2025 finals. Point values are quoted in the heading of
each solved sub-part. Partial-credit hints are inline. Refer back to \texttt{mock-exam-6.tex}
for the problem statements.}

\vspace{0.6em}
\hrule
\vspace{1em}

%=========================================================
\section*{Problem 1 \points{10} --- Fill-in-the-blank concepts}
%=========================================================

\sol{} \emph{(1\,P per blank.)}

\begin{enumerate}[(1)]
  \item \textbf{parametric}
  \item \textbf{prediction}
  \item \textbf{inference}
  \item \textbf{irreducible}
  \item \textbf{ridge}
  \item \textbf{nested cross-validation}
  \item \textbf{mini-batch stochastic gradient descent}
  \item \textbf{dropout}
  \item \textbf{early stopping}
  \item \textbf{label smoothing}
\end{enumerate}

\noindent\emph{Grading: 1\,P per correct blank. No partial credit for ``close'' alternatives.
For (5), ``lasso'' is the textbook trap --- lasso \emph{does} zero coefficients (corner of $L^1$
ball), ridge does not. For (6), ``stratified cross-validation'' is about preserving class balance,
not about honest hyperparameter assessment.}

\vspace{0.6em}
\hrule
\vspace{1em}

%=========================================================
\section*{Problem 2 \points{28} --- Multiple choice, T/F, short numeric}
%=========================================================

\subsection*{a) Bias--variance and benign overfitting \subpoints{3}}
\sol{} \emph{(0.75\,P per statement.)}
\begin{enumerate}[(i)]
  \item \textbf{True}. The decomposition is an algebraic identity. The first cross-term
        $2(f-\hat f)\cdot\varepsilon$ vanishes because $\mathbb{E}[\varepsilon]=0$ and
        $\varepsilon$ is independent of $\hat f$; the second cross-term
        $2(f-\mathbb{E}[\hat f])\cdot(\mathbb{E}[\hat f]-\hat f)$ vanishes after taking expectation
        over $\hat f$ because the second factor has mean zero by construction. Both rely on
        independence/zero-mean structure, not on any property of $\hat f$.
  \item \textbf{True}. In the classical view ridge introduces bias to lower variance, so one
        usually expects a trade. But \emph{moving from OLS to ridge changes the estimator}, which
        can shift both curves: if OLS suffers from collinearity its variance explodes, while a
        modestly regularized ridge fit can have both lower variance \emph{and} negligible bias
        (because the data don't really distinguish between the collinear coefficients anyway).
        The prof's ``trade-off'' skepticism leans on exactly this. Either example argument earns
        credit.
  \item \textbf{False}. This is benign overfitting / double descent. With $p\gg n$ and an
        appropriate optimizer (e.g.\ the min-norm interpolator picked by gradient descent), an
        interpolating model can still generalize.
  \item \textbf{True}. Larger irreducible noise $\sigma^2$ makes complex models spend variance
        chasing noise; the optimum on the bias/variance trade shifts toward simpler (less
        flexible) models.
\end{enumerate}

\subsection*{b) Cross-validation, with and without nesting \subpoints{4}}

\sol{} \subpoints{0.5} \textbf{(i)} \textbf{True}. Each LOO/large-$k$ fold uses $n-1$ ($\approx n$)
training points, so the CV-fitted models are closer to the model that would be trained on all $n$
data points $\Rightarrow$ less bias. Bias of CV decreases as $k$ grows.

\medskip
\sol{} \subpoints{0.5} \textbf{(ii)} \textbf{True}. The $n$ leave-one-out training sets overlap
almost completely, so their per-fold errors are highly correlated; averaging strongly correlated
numbers does not reduce variance much, and the LOOCV estimate has higher variance than the
$5$-fold one. (Prof's L11 line: \emph{``leave-one-out has better bias but K-fold will have a much
better variance, and typically you're winning by having less variance.''})

\medskip
\sol{} \subpoints{1} \textbf{(iii)} \textbf{False}. This is the canonical ``wrong-way CV''
pipeline: the predictor screening uses $y$, so the held-out folds have already informed which
predictors are retained. The CV estimate is biased \emph{downward} (it can be near zero even when
the truth is pure noise). Fix: redo the correlation screen \emph{inside} each training fold.

\medskip
\sol{} \subpoints{1} \textbf{(iv)} \textbf{1-SE rule}: Let $\hat\theta=\arg\min_\theta
\mathrm{CV}(\theta)$ achieving CV-MSE $m^\star$ with estimated standard error $\mathrm{SE}^\star$
across folds. Pick the \emph{simplest} (largest-$\lambda$, smallest model size) $\theta$ whose
CV-MSE is within $m^\star+\mathrm{SE}^\star$.

\medskip
\sol{} \subpoints{1} \textbf{(v)} Sum $=22.0$, so
$$\widehat{\mathrm{CV}\text{-MSE}}=\frac{22.0}{10}=\boxed{2.20}.$$
Sample variance of the $10$ fold MSEs (using $n-1=9$ in the denominator): squared deviations from
$2.20$ are $(0.10)^2,(0.20)^2,(0.20)^2,(0.40)^2,(0)^2,(0.30)^2,(0.10)^2,(0.30)^2,(0.20)^2,(0.20)^2$
$=0.01+0.04+0.04+0.16+0+0.09+0.01+0.09+0.04+0.04=0.52$. Sample SD $s=\sqrt{0.52/9}\approx 0.240$.
$$\widehat{\mathrm{SE}}(\widehat{\mathrm{CV}\text{-MSE}})=\frac{s}{\sqrt{10}}\approx\frac{0.240}{3.162}\approx\boxed{0.08}.$$
Accept anywhere in $[0.07, 0.09]$.

\subsection*{c) Mini-batch SGD and learning rate \subpoints{3}}
\sol{} \emph{(0.75\,P per statement.)}
\begin{enumerate}[(i)]
  \item \textbf{True}. Powers of two map cleanly onto GPU memory layouts; the choice is engineering
        rather than statistical. (Prof L23.)
  \item \textbf{False}. Larger batch sizes \emph{decrease} noise in the gradient estimate, but
        decreased noise means \emph{less} implicit $L^2$ regularization, not more. Small-batch SGD
        is the one that injects the noise that biases toward minimum-norm solutions. Direction
        reversed.
  \item \textbf{True}. Prof's L24 anecdote: $\eta=2$ ``horrible, bouncing everywhere,''
        $\eta\approx 0.1$ fine.
  \item \textbf{True}. Backprop \emph{is} the chain rule applied so each intermediate quantity is
        used (not recomputed). SGD is the update rule that consumes backprop's gradients.
\end{enumerate}

\subsection*{d) Odds, log-odds, interactions \subpoints{3}}

\sol{} \subpoints{1} \textbf{(i)} $p=\dfrac{\text{odds}}{1+\text{odds}}=\dfrac{0.25}{1.25}=\boxed{0.20}$.

\medskip
\sol{} \subpoints{1} \textbf{(ii)} Odds multiplier $=\exp(\hat\beta\cdot\Delta)=\exp(0.18\cdot 5)
=\exp(0.90)\approx \boxed{2.46}$.

\medskip
\sol{} \subpoints{1} \textbf{(iii)} For a treated patient, a unit increase in \texttt{age}
changes the log-odds by $\hat\beta_{\texttt{age}}+\hat\beta_{\texttt{age:treatment}}=0.04-0.03=0.01$,
so the odds multiply by $\exp(0.01)\approx \boxed{1.010}$.

\noindent\emph{Grading: full credit for $\approx 1.010$. Half credit for $\exp(0.04)\approx 1.041$
(ignoring the interaction --- the standard trap). Zero for treating the interaction in isolation.}

\subsection*{e) Collinearity and identifiability \subpoints{3}}
\sol{} \emph{(0.75\,P per statement.)}
\begin{enumerate}[(i)]
  \item \textbf{True}. Exact collinearity makes $X^\top X$ singular, so $(X^\top X)^{-1}$ does not
        exist and $\hat\beta=(X^\top X)^{-1}X^\top y$ is not uniquely defined. Any solution along
        the affine subspace $\beta_2 = (\text{const}) - 2\beta_1$ achieves the same fitted values.
  \item \textbf{True}. Under near-collinearity, the diagonal of $\sigma^2(X^\top X)^{-1}$ inflates
        (the individual SEs blow up), while the variance of an appropriate linear combination of
        the coefficients (essentially their sum) can stay controlled because the off-diagonal
        covariances are large and negative and ``cancel.'' This is the ``a million minus a million''
        intuition.
  \item \textbf{True}. With high correlation the individual $t$-tests can both fail to reject
        while the partial $F$-test on the pair rejects strongly. Concluding from the individual
        $t$-tests alone that neither predictor matters is exactly the trap; the variables matter
        \emph{jointly}.
  \item \textbf{False}. $K$ dummies together with an intercept produces a rank-deficient design
        matrix ($\sum$ of dummies $\equiv 1$, the intercept column). Use $K-1$ dummies and
        absorb the reference level into the intercept (or drop the intercept and use $K$ dummies).
\end{enumerate}

\subsection*{f) Bootstrap standard error \subpoints{3}}

\sol{} \subpoints{1} \textbf{(i)} With $\bar\theta^* = \frac{1}{B}\sum_{b=1}^B \hat\theta_b^*$,
$$\widehat{\mathrm{SE}}_{\mathrm{boot}}(\hat\theta) \;=\; \sqrt{\frac{1}{B-1}\sum_{b=1}^B
   \bigl(\hat\theta_b^* - \bar\theta^*\bigr)^2}.$$

\medskip
\sol{} \subpoints{1} \textbf{(ii)} Mean
$\bar\theta^* = \tfrac{1}{4}(2.1+2.5+1.8+2.4) = \tfrac{8.8}{4} = 2.20$.
Squared deviations from $2.20$: $(2.1-2.2)^2 = 0.01$,
$(2.5-2.2)^2 = 0.09$, $(1.8-2.2)^2 = 0.16$, $(2.4-2.2)^2 = 0.04$. Sum $= 0.30$.
Dividing by $B-1=3$ gives $0.30/3 = 0.10$, and
$$\widehat{\mathrm{SE}}_{\mathrm{boot}}(\hat\theta) = \sqrt{0.10} \approx \boxed{0.32}.$$

\medskip
\sol{} \subpoints{1} \textbf{(iii)} \textbf{(A) False} --- bootstrap is \emph{with} replacement.
\textbf{(B) True} --- the bootstrap quantifies the sampling \emph{variability} of $\hat\theta$
around itself, not its distance from the true $\theta$; biased estimators stay biased.
\textbf{(C) True} --- $B\in\{1000, 10000\}$ is the standard range.

\noindent\emph{Grading: 0.5 each for any one wrong, 1\,P all correct.}

\subsection*{g) Boosting: AdaBoost, gradient boosting, XGBoost \subpoints{3}}
\sol{} \emph{(0.75\,P per statement.)}
\begin{enumerate}[(i)]
  \item \textbf{True}. With $\alpha_m=\tfrac{1}{2}\log\tfrac{1-\mathrm{err}_m}{\mathrm{err}_m}>0$
        when $\mathrm{err}_m<0.5$, the weight-update $w_i\leftarrow w_i\exp(-\alpha_m\,y_i G_m(x_i))$
        multiplies misclassified weights by $e^{\alpha_m}>1$ and correctly classified weights by
        $e^{-\alpha_m}<1$; after renormalization the misclassified points end up with strictly
        higher \emph{relative} weight at the next iteration. (Some textbook conventions instead
        multiply only the misclassified weights and then renormalize; the relative-weight statement
        in the question is true under both.)
  \item \textbf{True}. The Hastie--Tibshirani--Friedman (HTF) gradient-boosting algorithm fits
        the new tree to the negative gradient $r_{im}=-[\partial L/\partial f]_{f=f_{m-1}}$ of the
        loss; for $L=\tfrac{1}{2}(y-f)^2$ this gradient is exactly $-(y-f)$, so the residual is the
        negative gradient up to sign.
  \item \textbf{False}. The directions are swapped. \emph{Random forest:} $B$ is not a tuning
        parameter (more trees can only help, by variance reduction); just pick ``enough.''
        \emph{Boosting:} $M$ \emph{is} a real tuning parameter because boosting \emph{can} overfit
        --- once residuals are noise, additional trees fit that noise.
  \item \textbf{True}. All three are core XGBoost ingredients; row/column subsampling, $L^2$
        leaf-weight penalty (and pruning parameter $\gamma$), and second-order Newton-style
        information are exactly what XGBoost adds on top of plain gradient boosting.
\end{enumerate}

\subsection*{h) Neural-network regularization \subpoints{3}}
\sol{} \emph{(0.75\,P per statement.)}
\begin{enumerate}[(i)]
  \item \textbf{False}. Dropout is a \emph{training-time-only} mechanism; at inference time the
        full network is used (with the surviving activations rescaled in the standard ``inverted
        dropout'' formulation).
  \item \textbf{True}. Prof L24: ``20\% is very common, never use 50\%.''
  \item \textbf{True}. This is exactly what early stopping does --- the validation-loss curve
        $U$-shapes and we return the weights at the minimum.
  \item \textbf{False}. Prof L26: \emph{``I can't think of an example where you'd ever want to
        train a neural network without a regularization.''} Even a small architecture relative to
        $n$ should use \emph{some} regularizer (mini-batch SGD itself counts).
\end{enumerate}

\subsection*{i) Principal component analysis \subpoints{3}}

\sol{} \subpoints{1} \textbf{(i)} Cumulative PVE: $0.60$, $0.85$, $0.95$, $1.00$. The first time
the cumulative crosses $0.90$ is at $\boxed{3}$ components.

\medskip
\sol{} \subpoints{1} \textbf{(ii)}
$z^*_1 = \phi_1^\top x^* = 0.60\cdot 1.0 + (-0.20)\cdot(-0.5) + 0.70\cdot 2.0 + 0.30\cdot 1.0
   = 0.60 + 0.10 + 1.40 + 0.30 = \boxed{2.40}$.

\medskip
\sol{} \subpoints{1} \textbf{(iii)} \textbf{(A) True} --- without standardization the largest-scale
predictor dominates the variance computation and tends to define PC1. \textbf{(B) True} --- loading
vectors are unit-norm by construction (otherwise PVE wouldn't be a proportion). \textbf{(C) False}
--- the first principal component is, by definition, the unit-norm linear combination of the
centered variables with the \emph{largest} sample variance; later components have progressively
smaller variances under the orthogonality constraint.

\noindent\emph{Grading for (iii): 0.33 per correct mark; full credit for the $(T,T,F)$ pattern.
Half credit if (C) is marked True --- a very common misread.}

\vspace{0.6em}
\hrule
\vspace{1em}

%=========================================================
\section*{Problem 3 \points{16} --- Theory, math, and pseudocode}
%=========================================================

\subsection*{a) The mathy one --- bias--variance decomposition and shrinkage \subpoints{8}}

\sol{} \subpoints{4} \textbf{(i) Derivation.} Write $y_0 = f(x_0) + \varepsilon_0$ with
$\mathbb{E}[\varepsilon_0]=0$ and $\mathrm{Var}(\varepsilon_0)=\sigma^2$. Add and subtract
$\mathbb{E}[\hat f(x_0)]$ inside the squared term:
$$(y_0 - \hat f(x_0))
   = \underbrace{\bigl(f(x_0) - \mathbb{E}[\hat f(x_0)]\bigr)}_{=:A}
   + \underbrace{\bigl(\mathbb{E}[\hat f(x_0)] - \hat f(x_0)\bigr)}_{=:B}
   + \underbrace{\varepsilon_0}_{=:C}.$$

\textbf{Square and take expectation.} $A$ is non-random (depends only on the true $f$ and on the
distribution of $\hat f$), so $\mathbb{E}[A^2] = A^2$. Expand:
\begin{align*}
\mathbb{E}[(y_0-\hat f(x_0))^2]
&= A^2 + \mathbb{E}[B^2] + \mathbb{E}[C^2] + 2A\,\mathbb{E}[B] + 2A\,\mathbb{E}[C] + 2\,\mathbb{E}[BC].
\end{align*}

\textbf{Cross-terms vanish.}
\begin{itemize}
  \item $\mathbb{E}[B] = \mathbb{E}[\hat f(x_0)] - \mathbb{E}[\hat f(x_0)] = 0$, so $2A\,\mathbb{E}[B]=0$.
  \item $\mathbb{E}[C] = \mathbb{E}[\varepsilon_0] = 0$, so $2A\,\mathbb{E}[C]=0$.
  \item For $\mathbb{E}[BC]$: $B$ is a function of $\hat f$ (so of the training set $\mathcal T$),
        $C=\varepsilon_0$ is independent of $\mathcal T$ by assumption, hence $B\perp C$ and
        $\mathbb{E}[BC]=\mathbb{E}[B]\mathbb{E}[C]=0$.
\end{itemize}

\textbf{The three surviving terms:}
\begin{align*}
A^2 &= \bigl(f(x_0)-\mathbb{E}[\hat f(x_0)]\bigr)^2 =: \mathrm{Bias}^2[\hat f(x_0)], \\
\mathbb{E}[B^2] &= \mathbb{E}\!\left[(\hat f(x_0)-\mathbb{E}[\hat f(x_0)])^2\right] = \mathrm{Var}(\hat f(x_0)), \\
\mathbb{E}[C^2] &= \mathrm{Var}(\varepsilon_0) = \sigma^2.
\end{align*}
Hence
$$\boxed{\;\mathbb{E}\!\left[(y_0-\hat f(x_0))^2\right]
   = \mathrm{Bias}^2[\hat f(x_0)] + \mathrm{Var}(\hat f(x_0)) + \sigma^2.\;}$$

\textbf{Sources of randomness.}
The outer expectation in $\mathbb{E}[(y_0-\hat f(x_0))^2]$ is taken \emph{jointly} over (a) the
random training sample $\mathcal T$ used to fit $\hat f$, and (b) the test-point noise
$\varepsilon_0$ in $y_0=f(x_0)+\varepsilon_0$. The inner $\mathbb{E}[\hat f(x_0)]$ and
$\mathrm{Var}(\hat f(x_0))$ are over the training-set distribution only. $f$ and $x_0$ are treated
as fixed.

\textbf{Naming the terms.} Squared bias and variance are both \emph{reducible} (one can lower
them by choosing a different estimator or different model class). The noise variance $\sigma^2$ is
\emph{irreducible} --- no estimator built from $X$ alone can predict $\varepsilon_0$.

\noindent\emph{Grading: 2\,P for both cross-terms vanishing with justification; 1\,P for the
``random over what'' explanation; 1\,P for the names and the reducibility comment. Deduct 0.5 if
$\sigma^2$ is dropped, 0.5 if the bias is reported without the square. Watch for sign errors in
the expansion.}

\medskip
\sol{} \subpoints{3} \textbf{(ii) Application to $\hat\mu_c = c\bar y$.}

We have $\bar y = \mu + \bar\varepsilon$ with $\bar\varepsilon\sim N(0,\sigma^2/n)$, so
$\mathbb{E}[\bar y]=\mu$ and $\mathrm{Var}(\bar y)=\sigma^2/n$. Hence:
\begin{align*}
\mathrm{Bias}(\hat\mu_c) &= \mathbb{E}[c\bar y] - \mu = c\mu - \mu = -(1-c)\mu,\\
\mathrm{Bias}^2(\hat\mu_c) &= (1-c)^2\,\mu^2,\\
\mathrm{Var}(\hat\mu_c) &= c^2\,\mathrm{Var}(\bar y) = \frac{c^2\sigma^2}{n}.
\end{align*}
Reducible MSE $= (1-c)^2\mu^2 + c^2\sigma^2/n$. Differentiate w.r.t.\ $c$ and set to zero:
$$-2(1-c)\mu^2 + 2c\,\sigma^2/n = 0 \;\Longrightarrow\; c\bigl(\mu^2 + \sigma^2/n\bigr) = \mu^2
   \;\Longrightarrow\; \boxed{\;c^\star = \frac{\mu^2}{\mu^2 + \sigma^2/n}\;}\in(0,1).$$
Sanity: when $\sigma^2/n\to 0$ (lots of data, low noise), $c^\star\to 1$ (no shrinkage, just use
the sample mean). When $\sigma^2/n \gg \mu^2$ (noisy, small sample), $c^\star\to 0$ --- the
optimal estimator collapses toward zero, sacrificing all bias to kill variance. \emph{Either
comment earns the closing point.}

\noindent\emph{Grading: 1\,P for $\mathrm{Bias}^2 = (1-c)^2\mu^2$; 1\,P for
$\mathrm{Var}=c^2\sigma^2/n$; 1\,P for the optimum $c^\star$. Accept stating $c^\star$ without the
limit comment if the algebra is right.}

\medskip
\sol{} \subpoints{1} \textbf{(iii) ``Trade-off'' skepticism.} The decomposition is exact, but it
is not always true that lowering one term \emph{requires} raising the other. Two cases the prof
flagged in the lectures:
\begin{enumerate}[(a)]
  \item \emph{Regularization can reduce both.} Adding a ridge penalty $\lambda\|\beta\|^2$ to a
        collinear OLS problem shrinks coefficients (small bias hit) but cuts the variance massively
        because $(X^\top X+\lambda I)^{-1}$ tames the inverse --- both terms can drop relative to
        plain OLS.
  \item \emph{Benign overfitting / double descent.} In the over-parameterized regime ($p\gg n$),
        once the model interpolates the training data, additional flexibility only changes which
        zero-training-error solution is picked (minimum-norm by the implicit bias of SGD);
        variance often \emph{decreases} again as $p$ grows.
\end{enumerate}
Either of the two earns the point.

\subsection*{b) Pseudocode for nested cross-validation \subpoints{5}}

\sol{} \subpoints{3} \textbf{(i) Pseudocode.}

\medskip
\noindent\fbox{\parbox{\dimexpr\textwidth-2\fboxsep-2\fboxrule}{%
\small
\textbf{Input:} data $D=\{(x_i,y_i)\}_{i=1}^n$; hyperparameter grid $\Lambda=\{\lambda_1,\dots,\lambda_T\}$; outer folds $K$, inner folds $K'$.

\textbf{Output:} an estimate of generalization error of the procedure
``ridge regression with $\lambda$ chosen by inner CV.''

\begin{enumerate}
  \item Randomly partition $D$ into $K$ disjoint \emph{outer} folds $F_1,\dots,F_K$.
  \item \textbf{For} $k=1,\dots,K$:
  \begin{enumerate}
    \item Let $D_k^{\mathrm{out}} = F_k$ (held-out test fold) and $D_k^{\mathrm{tr}} = D \setminus F_k$ (outer training).
    \item Randomly partition $D_k^{\mathrm{tr}}$ into $K'$ \emph{inner} folds $G_1,\dots,G_{K'}$.
    \item \textbf{For} each $\lambda\in\Lambda$:
       \textbf{For} $j=1,\dots,K'$: fit ridge with penalty $\lambda$ on $D_k^{\mathrm{tr}}\setminus G_j$; record MSE on $G_j$ as $\mathrm{MSE}_{k,j}(\lambda)$. \\
       Compute inner-CV score $\mathrm{IC}_k(\lambda) = \tfrac{1}{K'}\sum_{j=1}^{K'}\mathrm{MSE}_{k,j}(\lambda)$.
    \item Pick $\lambda^\star_k = \arg\min_{\lambda\in\Lambda} \mathrm{IC}_k(\lambda)$. \emph{(The choice can differ across outer folds.)}
    \item Refit ridge with penalty $\lambda^\star_k$ on all of $D_k^{\mathrm{tr}}$, predict on $D_k^{\mathrm{out}}=F_k$, and record outer-fold MSE $e_k=\mathrm{MSE}_{D_k^{\mathrm{out}}}$.
  \end{enumerate}
  \item Return $\widehat{\mathrm{Err}} = \tfrac{1}{K}\sum_{k=1}^K e_k$ as the nested-CV estimate of generalization error of the whole pipeline.
\end{enumerate}}}

\noindent\emph{Grading: 1\,P for explicit outer/inner loop structure; 1\,P for ``$\lambda^\star_k$ chosen using only $D_k^{\mathrm{tr}}$, and the outer fold $F_k$ never participates in selection''; 1\,P for the final aggregation across outer folds. Deduct 0.5 if the analyst's chosen $\lambda$ is presented as a single global $\hat\lambda$ instead of $\lambda^\star_k$ varying with $k$.}

\medskip
\sol{} \subpoints{1} \textbf{(ii)} If you pick $\lambda$ by single $K$-fold CV and \emph{report
the minimum CV-MSE} as the test-error estimate, you are taking the minimum of a noisy set of
estimates --- i.e.\ you are reporting an ``argmin'' statistic without correcting for the selection
bias. The reported CV-MSE will be \emph{optimistically biased downward}: the lower the noise at
the picked $\lambda$ relative to others, the more $\lambda$ ``benefits'' from the selection. Honest
assessment requires an outer loop whose held-out folds never participated in selecting $\lambda$.

\medskip
\sol{} \subpoints{1} \textbf{(iii)} If you screen predictors using their correlation with $y$
\emph{outside} the CV loop, the held-out CV folds have already informed the variable-selection
step, so the supposed ``held-out'' data is no longer held out --- it has helped shape the model.
The CV estimate is therefore biased \emph{downward}, and can be near zero even when the truth is
pure noise (the canonical ISLP exercise: $p=5000$ random predictors, two equiprobable labels;
naive pipeline reports ~0\% error). \textbf{Structural fix:} redo the correlation screen
\emph{inside each training fold}, using only that fold's training portion.

\subsection*{c) Backpropagation in a tiny network \subpoints{3}}

\sol{} \subpoints{2} \textbf{(i) Derivation.} Forward pass:
$z = w_1 x + b_1$, $h = \sigma(z)$, $\hat y = w_2 h + b_2$, $L = \tfrac{1}{2}(\hat y - y)^2$.

\textbf{Output-layer gradient.} Chain rule:
$$\frac{\partial L}{\partial w_2}
   = \frac{\partial L}{\partial \hat y}\cdot \frac{\partial \hat y}{\partial w_2}
   = (\hat y - y)\cdot h.$$

\textbf{Hidden-layer gradient.}
$$\frac{\partial L}{\partial w_1}
   = \frac{\partial L}{\partial \hat y}\cdot \frac{\partial \hat y}{\partial h}\cdot
     \frac{\partial h}{\partial z}\cdot \frac{\partial z}{\partial w_1}
   = (\hat y - y)\cdot w_2 \cdot h(1-h)\cdot x.$$

\noindent\emph{Grading: 1\,P for $\partial L/\partial w_2 = (\hat y - y)\cdot h$; 1\,P for the
full chain through to $\partial L/\partial w_1 = (\hat y - y)\cdot w_2 \cdot h(1-h)\cdot x$.
Accept either $\sigma'(z)$ or $h(1-h)$.}

\medskip
\sol{} \subpoints{1} \textbf{(ii) Numerical.}
$z = 0.5\cdot 2 + 0 = 1.0$.
$h = \sigma(1) = 1/(1+e^{-1}) = 1/(1+0.3679) = 1/1.3679 \approx 0.7311$.
$\hat y = 1\cdot 0.7311 + 0 \approx \boxed{0.731}$.
$\hat y - y \approx -0.269$.
$h(1-h) \approx 0.7311\cdot 0.2689 \approx 0.1966$.
$\frac{\partial L}{\partial w_1} = (-0.269)\cdot 1 \cdot 0.1966 \cdot 2 \approx \boxed{-0.106}$.

Accept any answer in $[-0.110, -0.100]$. Sign matters: a negative gradient means SGD will increase
$w_1$ to reduce loss, which is correct (current $\hat y$ undershoots $y=1$, and increasing $w_1$
pushes $z$ up, hence $h$ and $\hat y$ up).

\vspace{0.6em}
\hrule
\vspace{1em}

%=========================================================
\section*{Problem 4 \points{18} --- Concrete compressive strength}
%=========================================================

\subsection*{a) Linear regression with collinearity, interactions, encoding \subpoints{8}}

\sol{} \subpoints{1} \textbf{(i)} Count rows in the coefficient table: intercept, \texttt{cement},
\texttt{slag}, \texttt{water}, \texttt{superplast}, \texttt{coarse\_agg}, \texttt{fine\_agg},
$\log(\texttt{age})$, two \texttt{mix\_type} dummies, two \texttt{cement:mix\_type} interactions
$=\boxed{12}$ parameters. Sanity: residual d.f.\ $= 700-12 = 688$. \checkmark

\medskip
\sol{} \subpoints{2} \textbf{(ii) Collinearity diagnosis.}
\begin{itemize}
  \item \textbf{(a)} The variance-covariance matrix of OLS is $\mathrm{Var}(\hat\beta) = \sigma^2
        (X^\top X)^{-1}$. When two columns of $X$ are near-collinear, $X^\top X$ has a near-zero
        eigenvalue in their joint direction; the corresponding diagonal entries of $(X^\top X)^{-1}$
        blow up, so the individual SEs of $\hat\beta_{\texttt{water}}$ and
        $\hat\beta_{\texttt{superplast}}$ inflate, killing their individual $t$-statistics. The
        \emph{sum} $\beta_{\texttt{water}}+\beta_{\texttt{superplast}}$ (or any other well-conditioned
        combination orthogonal to the collinearity direction) remains well-estimated, which is
        why the joint $F$-test rejects: the model can be sure that \emph{some} combination of
        \texttt{water} and \texttt{superplast} matters, just not which of the two carries the
        effect.
  \item \textbf{(b)} \textbf{No.} Dropping both predictors on the basis of their individual
        $t$-tests would be wrong: the joint $F$-test on the pair rejects with $p<10^{-9}$, which
        means at least one of them belongs in the model. One sensible response is to drop one and
        keep the other (and live with the fact that you can't disentangle them), or to use a
        regularized fit (ridge / lasso / elastic net) that handles correlated predictors
        gracefully.
\end{itemize}

\noindent\emph{Grading: 1\,P for the $\mathrm{Var}(\hat\beta) = \sigma^2(X^\top X)^{-1}$ argument
plus the SE-inflation conclusion; 1\,P for explicitly rejecting the ``drop both'' inference and
giving a reasonable remedy.}

\medskip
\sol{} \subpoints{2} \textbf{(iii) Cement effect by mix type.}
A unit increase in \texttt{cement} changes the predicted \texttt{strength} (in MPa) by:
\begin{align*}
\text{Standard mix:}\quad &\hat\beta_{\texttt{cement}} = 0.110\;\text{MPa per kg/m}^3.\\
\text{High-strength:}\quad &\hat\beta_{\texttt{cement}} + \hat\beta_{\texttt{cement:mix\_high}}
   = 0.110 + 0.030 = \boxed{0.140}\;\text{MPa per kg/m}^3.\\
\text{Lightweight:}\quad &\hat\beta_{\texttt{cement}} + \hat\beta_{\texttt{cement:mix\_light}}
   = 0.110 + (-0.020) = \boxed{0.090}\;\text{MPa per kg/m}^3.
\end{align*}
The question asks for high-strength and lightweight; both shown.

\noindent\emph{Grading: 1\,P each. Half credit if the student reports $0.030$ or $-0.020$ alone
(the interaction term in isolation --- the standard trap). Zero if both terms are reported as
$0.110$ (interaction entirely ignored).}

\medskip
\sol{} \subpoints{1} \textbf{(iv) Encoding flip.} Switching the reference level from
\emph{standard} to \emph{lightweight}:
\begin{itemize}
  \item \textbf{(a) Residual standard error:} \textbf{unchanged} (the column space of $X$ is
        identical; only the basis changes).
  \item \textbf{(b) Coefficient on \texttt{cement}:} \textbf{changes} (it is now the slope
        \emph{for the lightweight reference mix}, which is $0.110-0.020 = 0.090$ instead of $0.110$).
  \item \textbf{(c) Coefficient on \texttt{mix\_high}:} \textbf{changes} (it is now the contrast
        ``high-strength vs.\ lightweight'' instead of ``high-strength vs.\ standard'').
  \item \textbf{(d) Predicted strength of any given mix:} \textbf{unchanged} --- fitted values are
        invariant under reparameterization.
\end{itemize}

\noindent\emph{Grading: 0.25\,P per part. All-or-nothing on each. The (a) and (d) ``unchanged''
items are the most commonly missed; flag if the student claims that ``fits change'' under
re-encoding.}

\medskip
\sol{} \subpoints{1} \textbf{(v) CI vs.\ PI.}
\begin{itemize}
  \item \textbf{(a)} A $95\%$ \emph{confidence interval} for $\mathbb{E}[\texttt{strength}\mid x_0]$
        contains the (population-level) mean strength of the sub-population of concrete mixes with
        predictor values $x_0$; a $95\%$ \emph{prediction interval} contains the future
        \texttt{strength} of \emph{one new individual} mix with predictor values $x_0$.
  \item \textbf{(b)} The prediction interval is wider. Its variance picks up an extra additive
        $\sigma^2$ term for the irreducible noise on the individual observation: ${\hat y}_0\pm t\,
        \hat\sigma\sqrt{1 + x_0^\top(X^\top X)^{-1}x_0}$ vs.\ ${\hat y}_0\pm t\,
        \hat\sigma\sqrt{x_0^\top(X^\top X)^{-1}x_0}$. As $n\to\infty$ the CI shrinks to a point but
        the PI stays bounded below by $\pm t\,\hat\sigma$.
\end{itemize}

\medskip
\sol{} \subpoints{1} \textbf{(vi) Fanning residuals.} The constant-variance (homoscedasticity)
assumption is violated. Common remedies: transform the response (e.g.\ $\log\texttt{strength}$ or
$\sqrt{\texttt{strength}}$), or use weighted least squares with weights inversely proportional to
the local residual variance.

\subsection*{b) Ridge regression with $10$-fold CV \subpoints{5}}

\sol{} \subpoints{1} \textbf{(i)} The ridge penalty $\lambda\sum_j\beta_j^2$ is scale-dependent;
without standardization a predictor in larger units (e.g.\ \texttt{cement} in kg/m$^3$ vs.\ a
binary dummy) is penalized far less because its coefficient is naturally small. Standardizing to
mean $0$ and variance $1$ puts every predictor on the same footing so shrinkage reflects signal,
not units.

\medskip
\sol{} \subpoints{1} \textbf{(ii)} Ridge minimizes
$$\hat\beta^R_\lambda \;=\; \arg\min_{\beta}\; \sum_{i=1}^n (y_i - \beta_0 - x_i^\top\beta)^2
   \;+\; \lambda\sum_{j=1}^p \beta_j^2.$$
As $\lambda\to 0$, the penalty disappears and $\hat\beta^R_\lambda\to\hat\beta^{\mathrm{OLS}}$.
As $\lambda\to\infty$, every coefficient is shrunk to $0$ and the fit collapses to the
unpenalized intercept ($\bar y$).

\medskip
\sol{} \subpoints{1} \textbf{(iii) 1-SE rule.} Among all $\lambda$ whose CV-MSE is within one
standard error (across folds) of the minimum CV-MSE, choose the \emph{largest} (most regularized,
simplest model). Formally,
$$\hat\lambda_{1\mathrm{SE}} = \max\Bigl\{\lambda:\;\mathrm{CV}(\lambda) \leq \mathrm{CV}(\hat\lambda_{\min}) + \widehat{\mathrm{SE}}(\mathrm{CV}(\hat\lambda_{\min}))\Bigr\}.$$

\medskip
\sol{} \subpoints{2} \textbf{(iv) Bias--variance interpretation.}
\begin{itemize}
  \item \textbf{(a)} OLS (test MSE $41.5$) is beaten by ridge at $\hat\lambda_{\min}$ (test MSE
        $39.8$). This suggests OLS is \emph{over-fitting} on this dataset --- variance is the
        binding term, and a moderate amount of bias buys a bigger reduction in variance, lowering
        test MSE. The concrete-strength data are a plausible case for this: with the strong
        collinearity between \texttt{water} and \texttt{superplast}, $X^\top X$ is poorly
        conditioned, OLS coefficients are unstable, and ridge specifically tames that instability.
  \item \textbf{(b)} The $1$-SE choice sits in between because it is \emph{more} regularized than
        $\hat\lambda_{\min}$. It picks up extra bias (every coefficient is shrunk further), giving
        up a little of the variance saving; the net effect is worse than $\hat\lambda_{\min}$ on
        test MSE here. We accept the $1$-SE penalty in exchange for a simpler, more stable model.
\end{itemize}

\subsection*{c) Gradient boosting \subpoints{5}}

\sol{} \subpoints{2} \textbf{(i) One iteration.}

Initialize $\hat f(x) \leftarrow \bar y$. Let $r_i = y_i - \hat f(x_i)$ be the current residuals.
At iteration $b+1$:
\begin{enumerate}
  \item Fit a regression tree $\hat g^{(b+1)}(x)$ of (small) interaction depth $d$ on the data
        $(x_i, r_i)$ --- i.e.\ \emph{the new tree is fit to the current residuals}.
  \item Update the ensemble: $\hat f(x) \leftarrow \hat f(x) + \nu \cdot \hat g^{(b+1)}(x)$.
  \item Update residuals: $r_i \leftarrow r_i - \nu \cdot \hat g^{(b+1)}(x_i)$, or equivalently
        recompute $r_i = y_i - \hat f(x_i)$ with the updated $\hat f$.
\end{enumerate}
After $M$ iterations the final predictor is $\hat f(x) = \bar y + \sum_{b=1}^{M} \nu\,\hat g^{(b)}(x)$.

\emph{Where $\nu$ enters: in step 2 (and equivalently step 3), shrinking each new tree's
contribution to the ensemble. With squared-error loss the residual is the negative gradient of
the loss, so this is a special case of the general gradient-boosting recipe.}

\noindent\emph{Grading: 1\,P for ``fit the new tree to current residuals''; 1\,P for the update
with the $\nu$ shrinkage shown explicitly. Accept either residuals-first or gradient-first framing.}

\medskip
\sol{} \subpoints{2} \textbf{(ii) Hyperparameters.}
\begin{itemize}
  \item $\nu$ \emph{(shrinkage / learning rate):} fix small ($\nu \le 0.1$ typical) --- not really
        tuned by CV, just chosen small.
  \item $d$ \emph{(interaction depth):} tune by CV. Small $d$ (e.g.\ $1$--$4$): each tree captures
        only low-order interactions, lowest variance per tree; larger $d$ admits higher-order
        interactions, more variance per tree.
  \item $M$ \emph{(number of trees):} tune by CV (or early stopping on a validation curve).
        Boosting \emph{can} overfit if $M$ is too large.
  \item \textbf{Coupling.} If $\nu$ is halved, each tree contributes half as much, so roughly
        \emph{twice as many} trees are needed to reach the same overall fit. Practically: $M\cdot\nu
        \approx$ constant for a given total ``capacity.''
\end{itemize}

\medskip
\sol{} \subpoints{1} \textbf{(iii) Ridge vs.\ boosting gap.} Ridge merely shrinks linear
coefficients; boosting fits a sum of (small) trees and can therefore capture \emph{non-linear
shapes} and \emph{interactions among predictors} that no linear model can. The large drop from
$39.8$ (ridge) to $26.3$ (boosting) suggests the predictors enter the truth non-linearly and/or
with non-trivial interactions --- which is consistent with the prior knowledge that, say,
\texttt{cement}, \texttt{water}, and \texttt{age} affect strength through complex chemical-curing
dynamics that no linear-in-the-coefficients model would capture.

\vspace{0.6em}
\hrule
\vspace{1em}

%=========================================================
\section*{Problem 5 \points{28} --- Telecom customer churn}
%=========================================================

\subsection*{a) Logistic regression with interaction \subpoints{8}}

\sol{} \subpoints{2} \textbf{(i)} For a one-unit increase in \texttt{monthly} the log-odds
change by the slope at that contract level.
\begin{align*}
\text{Month-to-month}\;(\text{reference}):\;& \Delta\hat\eta = 0.030;\\
\text{Two-year}:\;& \Delta\hat\eta = \hat\beta_{\texttt{monthly}} + \hat\beta_{\texttt{monthly:contract\_2yr}}
   = 0.030 + (-0.040) = -0.010.
\end{align*}
For an increase of $\Delta=10$ EUR:
\begin{align*}
\text{Month-to-month:}\;\text{odds}\times \exp(10\cdot 0.030) = \exp(0.30) \approx \boxed{1.350}.\\
\text{Two-year:}\;\text{odds}\times \exp(10\cdot (-0.010)) = \exp(-0.10) \approx \boxed{0.905}.
\end{align*}
\emph{Interpretation:} on a month-to-month contract, extra EUR is associated with \emph{more}
churn; on a two-year contract the sign flips and the same EUR increase is associated with
slightly \emph{lower} churn (the customer is locked in and modest price changes don't drive them
out).

\noindent\emph{Grading: 1\,P each. Half credit if the student reports $\exp(0.030)$ and
$\exp(-0.010)$ alone (per-unit factors, ignoring the requested $\Delta=10$). Zero if the
two-year case drops the interaction term entirely.}

\medskip
\sol{} \subpoints{1} \textbf{(ii)} The main-effect coefficient $0.030$ is the per-EUR effect of
\texttt{monthly} on log-odds \emph{only on the reference (month-to-month) contract}, because the
interactions $\texttt{monthly:contract\_1yr}$ and $\texttt{monthly:contract\_2yr}$ shift the
slope on the other levels. To describe the effect of \texttt{monthly} on churn ``on average''
across contracts you must combine $\hat\beta_{\texttt{monthly}}$ with the relevant interaction
coefficient on a per-contract basis:
$$\text{slope for contract level } c \;=\; \hat\beta_{\texttt{monthly}} + \hat\beta_{\texttt{monthly:}c}\quad\text{(zero for the reference level).}$$

\medskip
\sol{} \subpoints{3} \textbf{(iii) Predicted probability.} For the given customer, the linear
predictor is built term by term:
\begin{align*}
\hat\eta &= -1.90  &&\text{(intercept)} \\
& + (-0.040)\cdot 12 = -0.480 &&\text{(\texttt{tenure})}\\
& + 0.030\cdot 70 = 2.100 &&\text{(\texttt{monthly})}\\
& + (-0.010)\cdot 45 = -0.450 &&\text{(\texttt{age})}\\
& + 1.50 &&\text{(\texttt{contract\_1yr})}\\
& + (-0.020)\cdot 70 = -1.400 &&\text{(\texttt{monthly:contract\_1yr})}\\
& + 0 &&\text{(\texttt{senior} $=0$, \texttt{contract\_2yr} $=0$)}\\
&= \boxed{-0.630}.
\end{align*}
Sigmoid:
$$\hat p = \frac{1}{1+e^{-\hat\eta}} = \frac{1}{1+e^{0.630}} = \frac{1}{1+1.878} = \frac{1}{2.878} \approx \boxed{0.348}.$$
Accept any answer in $[0.34, 0.36]$.

\noindent\emph{Grading: 1\,P for assembling all relevant terms (incl.\ the one-year contract main
effect and its monthly-interaction; deduct if either is dropped); 1\,P for the value of
$\hat\eta\approx -0.63$; 1\,P for the correctly applied sigmoid. Deduct 0.5 for using
\texttt{contract\_2yr} instead of \texttt{contract\_1yr}.}

\medskip
\sol{} \subpoints{1} \textbf{(iv)} $\hat p \approx 0.35 < 0.5$, so the customer is predicted
\textbf{not} to churn at the default threshold. But the base rate of churn is $\approx 26\%$, so
$\hat p = 0.35$ is already well above the base rate; a more sensible classification threshold
might be lower (e.g.\ around $0.26$, the base rate), in which case the same customer would be
flagged as a likely churner --- $0.5$ is a poor default here.

\medskip
\sol{} \subpoints{1} \textbf{(v) Re-encoding.} If the colleague swaps the reference level of
\texttt{contract} from \emph{month-to-month} to \emph{two-year}:
\begin{itemize}
  \item \textbf{(a)} \textbf{No change.} The predicted probability of churn for any given customer
        is invariant under reparameterization of the design matrix.
  \item \textbf{(b)} \textbf{Yes, changes.} The two dummies are now ``month-to-month vs.\ two-year''
        and ``one-year vs.\ two-year,'' which is a different basis; the estimates and SEs are
        different numbers (though they encode the same fitted model).
  \item \textbf{(c)} \textbf{No change.} The overall fit (residual deviance, $-2\log L$, AIC) is
        invariant under reparameterization of the same model.
\end{itemize}

\subsection*{b) AdaBoost with pseudocode \subpoints{6}}

\sol{} \subpoints{3} \textbf{(i) Pseudocode.}

\medskip
\noindent\fbox{\parbox{\dimexpr\textwidth-2\fboxsep-2\fboxrule}{%
\small
\textbf{Input:} training data $\{(x_i, y_i)\}_{i=1}^n$ with $y_i\in\{-1,+1\}$; weak-learner
class $\mathcal G$; number of rounds $M$.

\textbf{Output:} ensemble classifier $G(x)$.

\begin{enumerate}
  \item Initialize weights $w_i^{(1)} = \tfrac{1}{n}$ for $i=1,\dots,n$.
  \item \textbf{For} $m=1,\dots,M$:
  \begin{enumerate}
    \item Fit weak classifier $G_m\in\mathcal G$ to the training data using weights $w_i^{(m)}$.
    \item Compute the weighted misclassification error
          $$\mathrm{err}_m \;=\; \frac{\sum_{i=1}^n w_i^{(m)}\,\mathbb{1}\!\{y_i\neq G_m(x_i)\}}
                                       {\sum_{i=1}^n w_i^{(m)}}.$$
    \item Compute the classifier weight
          $$\alpha_m \;=\; \tfrac{1}{2}\log\!\frac{1-\mathrm{err}_m}{\mathrm{err}_m}.$$
    \item Update observation weights
          $$w_i^{(m+1)} \;=\; w_i^{(m)} \cdot \exp\!\bigl(-\alpha_m\, y_i\, G_m(x_i)\bigr),
            \quad i=1,\dots,n,$$
          and renormalize so $\sum_i w_i^{(m+1)} = 1$.
  \end{enumerate}
  \item Return the final classifier $G(x) = \mathrm{sign}\!\left(\sum_{m=1}^{M} \alpha_m\, G_m(x)\right)$.
\end{enumerate}}}

\noindent\emph{Grading: 1\,P for the weight initialization, the $\mathrm{err}_m$ definition, and
the classifier weight $\alpha_m$ (the trio is standard); 1\,P for the multiplicative weight update
with $y_i G_m(x_i)\in\{\pm 1\}$; 1\,P for the final weighted-vote classifier with $\mathrm{sign}$.
Accept the equivalent ``factor-of-$2$'' convention $\alpha_m = \log\tfrac{1-\mathrm{err}_m}{\mathrm{err}_m}$
(without the $\tfrac{1}{2}$); the prof has used both. Deduct 0.5 if labels are coded as $\{0,1\}$
instead of $\{-1,+1\}$ --- the weight update only works with the signed convention.}

\medskip
\sol{} \subpoints{1} \textbf{(ii)} With $\mathrm{err}_m = 0.30$:
$$\alpha_m \;=\; \tfrac{1}{2}\log\frac{1-0.30}{0.30}
   = \tfrac{1}{2}\log\frac{0.70}{0.30}
   = \tfrac{1}{2}\log(2.333)
   = \tfrac{1}{2}\cdot 0.847
   \approx \boxed{0.42}.$$
Accept any answer in $[0.40, 0.44]$. \emph{If the student uses the unscaled convention $\alpha_m
=\log\tfrac{1-\mathrm{err}_m}{\mathrm{err}_m}\approx 0.85$, accept that too.}

\medskip
\sol{} \subpoints{1} \textbf{(iii) Weight updates.} Under the $\{-1,+1\}$ convention with
$\alpha_m\approx 0.42$:
\begin{itemize}
  \item \textbf{Misclassified} ($y_i G_m(x_i) = -1$): $\exp(+\alpha_m) = \exp(0.42)
        \approx \boxed{1.53}$.
  \item \textbf{Correctly classified} ($y_i G_m(x_i) = +1$): $\exp(-\alpha_m) = \exp(-0.42)
        \approx \boxed{0.65}$.
\end{itemize}
(With the unscaled $\alpha_m\approx 0.85$: misclassified factor $\approx 2.33$, correct factor
$\approx 0.43$. Either convention is accepted, provided consistent with the student's choice in
(ii).)

\medskip
\sol{} \subpoints{1} \textbf{(iv) Why weak learners (bias--variance).} Boosting builds a high-bias
low-variance \emph{ensemble} out of \emph{many} simple learners, each contributing a small,
shrunken correction. If the base learner were itself deep (low-bias, high-variance), each tree
would already fit much of the training signal --- including noise --- and only a few iterations
would suffice; the ensemble would inherit the per-tree variance and overfit aggressively. Shallow
``stumps'' have high bias individually but low variance; the sequential corrections drive bias
down without re-introducing the variance that deep trees would bring back in.

\subsection*{c) Neural-network classifier with regularization \subpoints{8}}

\sol{} \subpoints{2} \textbf{(i)} Each layer contributes $(\#\text{inputs} + 1)\cdot \#\text{outputs}$
parameters (the $+1$ counts the bias).
\begin{align*}
\text{input}\to\text{hidden}: & \quad (7+1)\cdot 16 = 128, \\
\text{hidden}\to\text{output}: & \quad (16+1)\cdot 1 = 17, \\
\text{Total:} & \quad \boxed{128 + 17 = 145} \text{ parameters.}
\end{align*}

\noindent\emph{Grading: 1\,P for the per-layer breakdown, 1\,P for the total. Deduct 0.5 if any
biases are dropped (e.g.\ a student forgetting all biases reports $7\cdot 16 + 16\cdot 1 = 128$,
or forgetting only the output bias reports $128+16=144$). The dropout, label smoothing, and early
stopping add \emph{no} parameters; the sigmoid on the output adds none.}

\medskip
\sol{} \subpoints{2} \textbf{(ii) Forward pass.}
$$z = b + \sum_{j=1}^{7} w_j x_j
   = -0.2 + 0.5\cdot 1 + (-0.3)\cdot 0 + 0.1\cdot 2 + 1.0\cdot 1 + (-0.5)\cdot 0 + 0.2\cdot(-1) + 0.4\cdot 1.$$
Term by term: $0.5 + 0 + 0.2 + 1.0 + 0 - 0.2 + 0.4 = 1.9$; with the bias, $z = -0.2 + 1.9 = \boxed{1.7}$.
ReLU: $h = \max(0, z) = \max(0, 1.7) = \boxed{1.7}$.

\medskip
\sol{} \subpoints{1} \textbf{(iii)} A network with only identity (linear) activations collapses
across layers: $\hat y = W_L\,W_{L-1}\cdots W_1\, x + (\text{biases}) = W_{\mathrm{eff}}\,x + b_{\mathrm{eff}}$,
where $W_{\mathrm{eff}}$ is a single matrix. The composite model is mathematically equivalent to
ordinary (multivariate) linear regression --- depth adds no representational capacity, because the
composition of linear maps is itself linear.

\medskip
\sol{} \subpoints{1} \textbf{(iv) The three regularizers.}
\begin{itemize}
  \item \textbf{Dropout:} During \emph{training only}, randomly zero out a fraction (here $20\%$)
        of the hidden-layer activations on each forward pass; the surviving activations are
        rescaled. At \emph{inference} the full network is used (no dropout, no rescaling beyond
        what was applied in training).
  \item \textbf{Early stopping:} During \emph{training}, monitor the loss on a held-out validation
        set; stop training when the validation loss has stopped improving, and return the weights
        from the validation-loss minimum (not the final iteration).
  \item \textbf{Label smoothing:} During \emph{training only}, replace the hard target $y\in\{0,1\}$
        by a softened target $\bigl((1-\varepsilon)\,y + \varepsilon/2\bigr)$ for each class (here
        $\varepsilon=0.05$); discourages the network from driving its logits to $\pm\infty$ to
        approach $\{0,1\}$.
\end{itemize}

\medskip
\sol{} \subpoints{1} \textbf{(v) Dropout rate of $50\%$.} Not a good idea --- the prof's convention
is ``$20\%$ is common, $50\%$ is too aggressive and almost never used.'' Dropping half of the
hidden units in every forward pass strips the network's representational capacity to the point
where it underfits.

\medskip
\sol{} \subpoints{1} \textbf{(vi) Training without regularization.} The classic over-fitting
signature: training loss continues to decrease monotonically (eventually toward zero) while the
validation loss drops initially, reaches a minimum, and then \emph{rises} as the model memorizes
training-set noise. The gap between train and validation loss widens with epochs.

\subsection*{d) LDA vs.\ QDA, plus random-forest variable importance \subpoints{4}}

\sol{} \subpoints{2} \textbf{(i)} LDA's two main modelling assumptions:
\begin{enumerate}
  \item \textbf{Gaussian class-conditional densities.} Within each class, the predictors are
        multivariate normal: $\mathbf{X}\mid Y=k \sim \mathcal{N}_p(\boldsymbol\mu_k, \boldsymbol\Sigma)$.
  \item \textbf{Common (pooled) covariance.} All classes share the \emph{same} $\boldsymbol\Sigma$.
\end{enumerate}
\textbf{Why prefer LDA when one class is small.} QDA estimates a separate covariance matrix
$\boldsymbol\Sigma_k$ per class --- here $K=2$ with the minority class ``churners'' having
roughly $900$ training points. QDA must estimate $p(p+1)/2$ covariance parameters from those $900$
points; with $p=5$ that is $15$ parameters from $900$ obs, not catastrophic, but the variance of
$\widehat{\boldsymbol\Sigma}_{\text{churn}}$ is much higher than the pooled $\widehat{\boldsymbol\Sigma}$ used
by LDA. The bias QDA pays for assuming distinct covariances is small if they really are similar;
the variance saving of LDA's pooled estimate is real either way. In a small-class regime LDA
typically wins on test error --- as is consistent here ($0.196$ for LDA, $0.221$ for QDA).

\medskip
\sol{} \subpoints{1} \textbf{(ii)} For classification the standard default is $\mathtt{mtry}
=\sqrt{p}=\sqrt{7}\approx 2.65$, conventionally rounded up to $\mathtt{mtry}=3$. We want
$\mathtt{mtry}<p$ so the individual trees are forced to use different predictors at the top splits
and become \emph{decorrelated}; averaging decorrelated trees yields better variance reduction than
averaging the highly-correlated trees that bagging ($\mathtt{mtry}=p$) would produce.

\medskip
\sol{} \subpoints{1} \textbf{(iii) What variable importance \emph{cannot} tell you.} It does not
tell you the \emph{direction} (sign) or \emph{shape} of any predictor's effect on $\hat p(\text{churn})$.
For example, the plot tells you that \texttt{tenure} matters a lot, but it does not tell you
whether long tenure raises or lowers the churn probability, nor what the functional form of that
relationship looks like (linear? saturating?). To recover those, you would compute partial
dependence plots (or ICE curves) or compare with a parametric model. \emph{Also acceptable: it
says nothing about causality --- a predictor that proxies for a confounder will have high
importance without having any causal effect.}

\subsection*{e) Class imbalance \subpoints{2}}

\sol{} \subpoints{1} \textbf{(i)} A constant ``no churn'' classifier mis-classifies exactly the
$26\%$ of customers who churn and correctly classifies the rest, so its test error rate is
$\approx \boxed{0.26}$.

\medskip
\sol{} \subpoints{1} \textbf{(ii)} The naive baseline ($0.26$) is \emph{worse} than LDA ($0.196$),
QDA ($0.221$), and the random forest ($0.179$), so plain accuracy does usefully discriminate
here --- but only modestly, and at the price of saying nothing about the model's ability to catch
the (commercially valuable) churners. In this kind of imbalanced setting one should privilege the
\emph{pair} (\textbf{sensitivity, specificity}), or a threshold-independent summary such as
\textbf{AUC}, or a cost-weighted metric reflecting that one retained customer is worth far more
than the cost of one unnecessary retention offer. Reporting only the overall test error rate
hides the fact that, at threshold $0.5$, all three classifiers are still missing a substantial
fraction of churners (sensitivity well below $1$).

\vspace{1em}
\hrule
\vspace{0.5em}

\noindent\textbf{End of solution proposal.}\quad Total awarded: $10+28+16+18+28=100$ points.

\end{document}
